% TLDR: Abstract distinguishes a prospective protocol from historical evidence.
\begin{abstract}
% Context and problem.
Artificial intelligence (AI) judges are useful evaluation instruments only to the extent that their ratings measure the intended property.
For formal specification coverage, a high correlation with experts does not establish agreement about score levels or disagreement between experts on individual artifacts.
% Proposed approach.
We specify a study in which a judge predicts the six-category human-rating distribution conditional on Python source, a candidate Lean formalization, and a gold reference.
The primary objective is normalized ranked probability score, with correlations and absolute error retained as diagnostics.
% Concrete design and evidence status.
A prospective design covers 896 benchmark tasks and four solver configurations, with 24 distinct experts and a starting budget of 14,848 individual ratings.
Separate task, solver, and expert holdouts test different transfer claims, while controlled specification changes test whether human agreement reflects the intended construct.
Historical VeriBench evidence motivates the design: a 30-artifact, ten-task test has Spearman correlation 0.826, but its 95\% task-bootstrap interval is [0.316, 0.966] and its normalized judge-minus-expert mean bias is +0.181.
% Scope and limitation.
Those previously exposed data do not validate the proposed judge; fresh human collection and the prospective results remain pending.
The contribution is a reproducible measurement protocol that separates predictive agreement, semantic validity, and uncertainty about both.
\end{abstract}
